% ═══════════════════════════════════════════════════════════════════
%  Chapitre 1 — Présentation de l'organisme d'accueil et du projet
% ═══════════════════════════════════════════════════════════════════
\chapter{Présentation de l'organisme d'accueil et du contexte du projet}

\section*{Introduction}
\addcontentsline{toc}{section}{Introduction}

Ce chapitre présente le cadre dans lequel ce projet de fin d'études a été
réalisé. Il décrit d'abord l'organisme d'accueil --- son activité, son
organisation et les personnes qui y travaillent --- puis le contexte du projet
proprement dit : la façon dont l'institut fonctionnait avant l'intervention,
les difficultés que ce fonctionnement engendrait, la solution retenue pour y
répondre, ainsi que les enjeux et les contraintes qui l'ont encadrée. Il se
termine par la méthode de conduite de projet adoptée et par la planification
qui en découle.

% ─────────────────────────────────────────────────────────────────
\section{Présentation de l'organisme d'accueil}

\subsection{Le centre \yani}

\yani\ est un centre de beauté et de bien-être installé à Marrakech. C'est une
structure de proximité : une clientèle essentiellement féminine, largement
fidélisée, qui revient à intervalles réguliers et qui connaît personnellement
la gérante. Cette proximité est le principal atout commercial de l'institut ;
c'est aussi ce qui rend la question de l'outil informatique délicate, car toute
solution qui interposerait une machine froide entre la cliente et l'institut
détruirait précisément ce qui fait sa valeur.

\figProjet[0.32]{logo.png}{Identité visuelle de \yani}{logo}

L'identité visuelle du centre --- un or lumineux sur fond crème ou noir profond
--- a été reprise à l'identique dans les trois applications développées. Ce
n'est pas un détail cosmétique : pour une cliente, l'application doit
manifestement être celle de l'institut qu'elle connaît, et non un outil
générique de réservation dans lequel l'institut ne serait qu'une entrée parmi
d'autres.

\subsection{Activités et offre de services}

L'activité du centre se répartit en deux volets, qui ont chacun leur logique
propre et qui sont tous deux couverts par la solution développée.

\begin{itemize}[leftmargin=1.2cm]
  \item \textbf{Les prestations de soin.} Elles constituent le c\oe ur de
  métier. Chaque prestation appartient à une catégorie (soins du visage, soins
  du corps, épilation, coiffure, onglerie\ldots), porte un tarif et
  s'effectue en cabine, sur rendez-vous. La ressource limitante est le nombre
  de cabines réservables simultanément.

  \item \textbf{La vente de produits cosmétiques.} L'institut revend des
  produits de soin. Cette activité obéit à une logique différente : elle porte
  sur un stock physique, se règle à la remise ou à la livraison, et n'occupe
  aucune cabine.
\end{itemize}

À ces deux volets s'ajoute un \textbf{programme de fidélité}, transversal, qui
récompense à la fois les prestations réalisées et les produits achetés.

\subsection{Organisation et acteurs}

L'institut est une structure de petite taille, ce qui a une conséquence directe
sur la conception : les rôles ne sont pas cloisonnés comme ils le seraient dans
une entreprise plus grande, et une même personne assure souvent plusieurs
fonctions dans la même journée. Trois profils ont néanmoins été distingués,
parce qu'ils correspondent à trois niveaux de responsabilité réellement
différents.

\begin{table}[H]
\centering
\caption{Les acteurs de l'institut et leurs responsabilités}
\label{tab:acteurs-institut}
\begin{tabularx}{\textwidth}{L{3.1cm} L{2.4cm} X}
\toprule
\entete{Acteur} & \entete{Rôle système} & \entete{Responsabilités} \\
\midrule
La gérante & \texttt{ADMIN} &
Direction du centre. Fixe les tarifs et le catalogue, règle les horaires
d'ouverture et la capacité, définit le programme de fidélité, gère les comptes
du personnel et consulte les exports comptables. \\

Le personnel & \texttt{STAFF} &
Réalise les soins et tient le comptoir. Traite les commandes, confirme et
clôture les rendez-vous, honore les bons de récompense, crédite
exceptionnellement des points à la main. \\

Les clientes & \texttt{CLIENT} &
Consultent les catalogues, réservent leurs rendez-vous, commandent des
produits et suivent leur compte de fidélité. \\
\bottomrule
\end{tabularx}
\end{table}

\figProjet[0.75]{organisation-centre.png}{Organisation du centre et acteurs du système}{organisation}

Cette hiérarchie de rôles n'est pas une simple convention documentaire : elle
est appliquée par le serveur sur chaque route de l'API, indépendamment de ce
que l'interface propose ou masque. Une cliente qui obtiendrait par un moyen
quelconque l'adresse d'une route d'administration se verrait refuser l'accès
par le serveur lui-même. Ce point est développé au chapitre~4.

% ─────────────────────────────────────────────────────────────────
\section{Présentation du contexte du projet}

\subsection{Vue d'ensemble du projet}

Le projet consiste à concevoir et à réaliser une solution logicielle complète
couvrant l'ensemble de l'activité de l'institut. Elle prend la forme d'un
mono-dépôt réunissant trois applications qui partagent une seule base de
données :

\begin{itemize}[leftmargin=1.2cm]
  \item une \textbf{API REST} qui porte les données et la totalité des règles
  métier ;
  \item une \textbf{application mobile} destinée aux clientes ;
  \item un \textbf{back-office web} destiné à la gérante et à son personnel.
\end{itemize}

Le choix de faire porter toutes les règles métier par l'API, et par elle seule,
est structurant. Les deux interfaces --- mobile et web --- n'implémentent
aucune règle : elles affichent ce que le serveur leur donne et lui soumettent
ce que l'utilisatrice saisit. Une règle telle que « on ne peut pas réserver
dans le passé » ou « un rendez-vous annulé ne peut plus être marqué terminé »
n'existe qu'à un seul endroit. Ce principe évite la divergence classique où
l'application mobile et le back-office finissent par appliquer des règles
subtilement différentes à la même donnée.

\subsection{Étude de l'existant}

Avant tout développement, une observation du fonctionnement réel de l'institut
a été menée. Elle a mis en évidence quatre processus manuels, tous fonctionnels
au quotidien, et tous porteurs des mêmes limites.

\paragraph{La prise de rendez-vous.} Elle s'effectue par trois canaux
concurrents : le téléphone, la messagerie instantanée et la présentation
directe au comptoir. La réservation est inscrite au crayon dans un carnet
papier, unique et physiquement situé à l'accueil. La cliente ne dispose
d'aucune confirmation écrite autre que la mémoire de l'échange.

\paragraph{La vente de produits.} Elle se fait exclusivement sur place. Le
stock n'est suivi par aucun outil : la gérante l'estime de visu. Un produit
épuisé ne se découvre qu'au moment où une cliente le demande.

\paragraph{La fidélité.} Elle repose sur une carte cartonnée tamponnée à chaque
visite. La carte perdue emporte l'historique avec elle ; la carte oubliée à la
maison prive la cliente de son tampon ; et l'institut n'a aucune vision
consolidée de ce que le programme lui coûte ni de ce qu'il lui rapporte.

\paragraph{Le suivi de l'activité.} Il n'existe pas de manière systématique. Il
n'est pas possible de répondre, sans reconstitution manuelle, à des questions
aussi élémentaires que « combien de rendez-vous ce mois-ci ? », « quelles sont
les prestations les plus demandées ? » ou « quelles clientes ne sont pas
revenues depuis six mois ? ».

\subsection{Problématique}

Ce fonctionnement, viable tant que l'activité reste modeste, engendre des
difficultés qui se renforcent mutuellement à mesure qu'elle croît.

\begin{itemize}[leftmargin=1.2cm]
  \item \textbf{Le risque de double réservation.} Trois canaux alimentent un
  carnet unique. Une réservation prise par téléphone pendant qu'une cliente
  écrit sur la messagerie peut viser le même créneau. Le conflit ne se découvre
  qu'au moment où les deux clientes se présentent, et sa résolution se fait
  toujours au détriment de l'une d'elles.

  \item \textbf{L'indisponibilité hors des heures d'ouverture.} Une cliente qui
  souhaite réserver le soir ou le dimanche doit attendre. Une partie de ces
  demandes ne se transforme jamais en rendez-vous.

  \item \textbf{Le temps consacré au téléphone.} Chaque appel interrompt un
  soin en cours. Ce coût est invisible dans les comptes, mais il est réel : il
  se paie en qualité de service pour la cliente présente en cabine.

  \item \textbf{L'absence de mémoire exploitable.} Le carnet et la carte de
  fidélité ne conservent qu'une trace éphémère et non consolidée. L'institut ne
  peut donc ni mesurer son activité, ni identifier ses prestations phares, ni
  détecter l'érosion de sa clientèle.

  \item \textbf{La fragilité du support physique.} Un carnet se perd, s'abîme,
  et n'existe qu'en un seul exemplaire, à un seul endroit.

  \item \textbf{La rupture de stock invisible.} Sans suivi, un produit épuisé
  n'est découvert qu'au comptoir, devant la cliente.
\end{itemize}

Ces difficultés ont un dénominateur commun : \emph{l'information de l'institut
n'existe qu'en un seul exemplaire, dans un seul lieu, et pendant les seules
heures d'ouverture}. C'est cette contrainte que le projet lève.

\subsection{Solution proposée}

La solution retenue est une plate-forme composée de trois applications adossées
à une base de données unique. Ce choix mérite d'être justifié : il aurait été
possible de se contenter d'un back-office web, ou de recourir à une plate-forme
de réservation existante. Ni l'une ni l'autre de ces options n'a été retenue.

Un back-office seul aurait déplacé le carnet papier vers un écran, sans rien
changer pour la cliente : elle aurait continué d'appeler. Une plate-forme de
réservation généraliste, quant à elle, aurait réglé la question des créneaux
mais placé l'institut au milieu d'un annuaire de concurrents, sans couvrir ni
la vente de produits ni un programme de fidélité propre --- et en faisant
dépendre l'activité d'un tiers sur lequel l'institut n'a aucune prise.

La plate-forme développée comprend donc :

\begin{itemize}[leftmargin=1.2cm]
  \item \textbf{Une API REST (NestJS, Prisma, PostgreSQL).} Elle expose plus de
  quatre-vingts routes et porte l'ensemble des règles métier : catalogue,
  moteur de créneaux, machine à états des rendez-vous et des commandes,
  programme de fidélité, authentification, exports.

  \item \textbf{Une application mobile (React Native, Expo).} Vingt-quatre
  écrans permettent à la cliente de parcourir les catalogues, de réserver, de
  commander, de suivre ses points et de gérer son compte. Elle est disponible
  en français, en arabe et en anglais.

  \item \textbf{Un back-office web (React, Vite).} Huit pages donnent à
  l'institut la maîtrise complète de son activité : tableau de bord,
  commandes, rendez-vous, catalogue, fidélité, horaires, utilisateurs, avec
  export Excel des tableaux réservé à la gérante.
\end{itemize}

\begin{table}[H]
\centering
\caption{Réponse apportée à chaque difficulté identifiée}
\label{tab:probleme-solution}
\begin{tabularx}{\textwidth}{L{5.2cm} X}
\toprule
\entete{Difficulté} & \entete{Réponse apportée} \\
\midrule
Double réservation d'un créneau &
Canal de réservation unique, et sérialisation des réservations concurrentes par
verrou consultatif PostgreSQL : deux demandes simultanées sur le même créneau
sont traitées l'une après l'autre, jamais en parallèle. \\

Indisponibilité hors ouverture &
Réservation possible à toute heure depuis l'application, sur la grille de
créneaux calculée par le serveur à partir des horaires réels du centre. \\

Temps passé au téléphone &
La réservation en autonomie décharge l'accueil ; le rendez-vous pris par
téléphone reste possible, le personnel le saisissant pour la cliente. \\

Absence de mémoire &
Historique complet en base, tableau de bord, et exports Excel horodatés des
clientes, des commandes et des rendez-vous. \\

Fragilité du support papier &
Base PostgreSQL, script de sauvegarde horodatée avec rotation à quatorze jours
et procédure de restauration documentée. \\

Rupture de stock invisible &
Stock suivi par produit, contrôlé à la commande, décrémenté de façon atomique à
la confirmation, et alerte de stock faible sur le tableau de bord. \\

Carte de fidélité perdue &
Compte de fidélité attaché au compte de la cliente : solde, historique, paliers
de visites et bons de récompense, consultables depuis l'application. \\
\bottomrule
\end{tabularx}
\end{table}

\subsection{Enjeux du projet}

\begin{itemize}[leftmargin=1.2cm]
  \item \textbf{Enjeu commercial.} Rendre la réservation possible en dehors des
  heures d'ouverture, et transformer en rendez-vous des demandes aujourd'hui
  perdues faute d'interlocuteur disponible.

  \item \textbf{Enjeu de fidélisation.} Remplacer une carte cartonnée fragile
  par un compte consultable, doté de paliers de visites et de bons de
  récompense traçables --- un programme que l'institut peut enfin piloter et
  chiffrer.

  \item \textbf{Enjeu de pilotage.} Donner à la gérante une visibilité qu'elle
  n'a jamais eue sur son activité : commandes en attente, rendez-vous du jour,
  stocks faibles, historique exportable.

  \item \textbf{Enjeu d'image.} Une application aux couleurs de l'institut,
  soignée et trilingue, participe directement à la perception de la marque
  auprès d'une clientèle jeune et connectée.

  \item \textbf{Enjeu de conformité.} Traiter des données personnelles impose
  des obligations, notamment le droit à l'effacement prévu par la loi 09-08 et
  contrôlé par la CNDP. Ce droit devait être réellement exerçable, et non
  simplement annoncé.
\end{itemize}

\subsection{Contraintes}

\paragraph{Contrainte d'exactitude sur les créneaux.} Un créneau de cabine ne
peut pas être vendu deux fois. Cette contrainte, banale à énoncer, est la plus
difficile du projet : la vérification de disponibilité et l'insertion du
rendez-vous portent sur un \emph{ensemble} de lignes, ce qu'aucune écriture
conditionnelle sur une ligne unique ne sait rendre atomique.

\paragraph{Contrainte de fuseau horaire.} L'institut vit à l'heure de
Marrakech, dont le décalage varie dans l'année. Les instants sont stockés en
UTC, les horaires d'ouverture exprimés en heure locale. Toute confusion entre
les deux se traduit par une grille de créneaux décalée d'une heure --- une
erreur silencieuse, qui ne lève aucune exception et qui ne se voit qu'en
comparant l'écran au carnet.

À noter, car cela surprend à la lecture du code : le fuseau est déclaré sous le
nom \texttt{Africa/Casablanca}. Ce n'est pas une erreur ni un reliquat. La base
IANA nomme chaque zone d'après sa ville la plus peuplée, et cette zone couvre
l'ensemble du Maroc, Marrakech comprise. La renommer serait la casser.

\paragraph{Contrainte de conservation de l'historique.} Une cliente peut
demander la suppression de son compte. Ses commandes et ses rendez-vous
portent pourtant le chiffre d'affaires de l'institut, qui doit lui survivre.
Concilier les deux exige une anonymisation, et non une suppression en cascade.

\paragraph{Contrainte d'usage.} L'utilisatrice principale du back-office n'est
pas informaticienne. L'outil devait être immédiatement compréhensible, sans
formation ni documentation à consulter --- ce qui a proscrit toute exposition
de valeurs techniques brutes dans l'interface comme dans les exports.

\paragraph{Contrainte de coût.} L'institut est une petite structure.
L'ensemble de la solution repose donc sur des technologies libres et sur une
infrastructure minimale, déployable sur un seul serveur.

\paragraph{Contrainte de délai.} Le projet devait aboutir à une plate-forme
fonctionnelle et déployable dans la durée impartie au stage de fin d'études,
soit environ huit semaines de développement effectif.

% ─────────────────────────────────────────────────────────────────
\section{Conduite du projet}

\subsection{La méthode agile SCRUM}

Le choix d'une méthode de développement conditionne l'organisation de tout le
projet. La méthode agile SCRUM a été retenue, pour une raison qui tient au
contexte plus qu'à la mode : les besoins de l'institut n'étaient pas
formalisables à l'avance. Ils se sont précisés à mesure que des versions
concrètes étaient montrées à la gérante.

L'exemple le plus net concerne le moteur de créneaux. La première version
calculait les disponibilités à partir de la durée de chaque prestation : une
prestation d'une heure vingt-cinq occupait le temps correspondant. Le
comportement était techniquement juste, et fonctionnellement inutilisable ---
il produisait des heures de rendez-vous irrégulières, impossibles à annoncer au
téléphone. La gérante espace en réalité ses rendez-vous d'un écart fixe, quelle
que soit la prestation. Le moteur a donc été refait : un rendez-vous occupe
exactement un créneau, et la durée de la prestation n'est plus qu'une
information de gestion. Aucune spécification écrite à l'avance n'aurait produit
cette règle ; seule la confrontation d'une version fonctionnelle au métier réel
l'a fait apparaître.

\figProjet[0.72]{scrum.png}{Le cycle de la méthode agile SCRUM}{scrum}

\subsection{Répartition des rôles}

Le projet ayant été mené par un seul développeur, la répartition canonique des
rôles SCRUM a été adaptée sans être abandonnée --- car ce sont bien trois
responsabilités distinctes, même lorsque deux d'entre elles sont portées par la
même personne.

\begin{table}[H]
\centering
\caption{Répartition des rôles SCRUM sur le projet}
\label{tab:roles-scrum}
\begin{tabularx}{\textwidth}{L{3.4cm} L{3.4cm} X}
\toprule
\entete{Rôle} & \entete{Tenu par} & \entete{Responsabilité} \\
\midrule
\textit{Product Owner} & Madame Fati, gérante &
Exprime le besoin, arbitre les priorités, valide chaque incrément à l'aune du
fonctionnement réel de l'institut. \\

\textit{SCRUM Master} & \nouvelleEncadrantPedagogique &
Cadence les revues, lève les obstacles et veille au respect de la méthode. \\

\textit{Équipe de développement} & \nouvelleAuteur &
Conception, développement des trois applications, tests, documentation et
déploiement. \\
\bottomrule
\end{tabularx}
\end{table}

\subsection{Réunions}

Trois rendez-vous ont rythmé chaque sprint.

\begin{itemize}[leftmargin=1.2cm]
  \item \textbf{La planification de sprint.} En début de sprint, les éléments
  les plus prioritaires du \textit{product backlog} sont sélectionnés pour
  constituer le \textit{sprint backlog}.

  \item \textbf{La revue de sprint.} En fin de sprint, l'incrément est
  démontré : l'application est installée sur le téléphone de la gérante et le
  back-office ouvert devant elle. C'est là que les écarts entre le besoin
  exprimé et le besoin réel se révèlent.

  \item \textbf{La rétrospective.} Immédiatement après la revue, elle porte sur
  la façon de travailler et non sur le produit. C'est de l'une d'elles qu'est
  née la règle appliquée dans tout le projet : \emph{tout comportement non
  évident doit être expliqué dans le code, à l'endroit où il s'applique}. Cette
  règle explique la densité inhabituelle de commentaires du dépôt, et elle a
  été décisive lors des reprises après interruption.
\end{itemize}

\subsection{Le développement piloté par les tests}

Une approche inspirée du TDD a été appliquée, non pas uniformément, mais là où
elle apporte le plus : les règles métier délicates, et tout particulièrement
les comportements en situation de concurrence.

Écrire d'abord le test qui prouve que trois réservations simultanées ne peuvent
pas occuper deux cabines oblige à formuler la propriété avant de choisir la
technique. C'est ce test, écrit contre une véritable base PostgreSQL, qui a
imposé le recours à un verrou consultatif : la solution intuitive --- compter
puis insérer --- le faisait échouer de façon reproductible.

\figProjet[0.62]{tdd.png}{Le cycle du développement piloté par les tests}{tdd}

Le même cycle a été appliqué au crédit de fidélité, à l'émission des bons de
récompense, à la rotation des jetons de session et à l'anonymisation des
comptes. Ces cinq domaines concentrent l'essentiel des huit suites de tests
qui écrivent réellement en base.

\subsection{Les sprints}

Le projet s'est déroulé sur six sprints, du 6 juillet au 27 août 2026.

\begin{table}[H]
\centering
\caption{Découpage du projet en sprints}
\label{tab:sprints}
\begin{tabularx}{\textwidth}{L{1.5cm} L{2.9cm} X}
\toprule
\entete{Sprint} & \entete{Période} & \entete{Contenu livré} \\
\midrule
1 & 6 -- 9 juillet &
Socle de l'API : authentification JWT, catalogues de prestations et de
produits, module de rendez-vous, horaires d'ouverture, fidélité, récompenses,
gestion des utilisateurs, jetons de rafraîchissement, gestion du fuseau du
centre, documentation Swagger. \\

2 & 12 -- 14 juillet &
Application mobile cliente : catalogues, réservation, fidélité, profil,
rendez-vous. Extraction des composants réutilisables et refonte graphique
complète aux couleurs de l'institut. \\

3 & 16 -- 23 juillet &
Commandes de produits avec paiement à la remise, et back-office complet :
tableau de bord, commandes, rendez-vous, catalogue, fidélité, utilisateurs. \\

4 & 7 -- 11 août &
Durcissement : vérification de l'adresse e-mail, réinitialisation du mot de
passe, récompenses de palier, écritures atomiques sur le solde, le stock et
les créneaux, anonymisation des comptes supprimés, identité native de
l'application mobile. \\

5 & 12 -- 18 août &
Réglages du centre et fermetures exceptionnelles, contrôle du stock,
pagination des listes, finitions d'expérience utilisateur, exports Excel
réservés à la gérante, internationalisation en français, arabe et anglais. \\

6 & 21 -- 27 août &
Hébergement local des images du catalogue, optimisation du décodage des
images, mise à jour du script de sauvegarde, conteneurisation et finalisation
de la documentation d'exploitation. \\
\bottomrule
\end{tabularx}
\end{table}

\subsection{Planification du projet}

La planification a été établie avant le démarrage puis ajustée à chaque revue.
Le diagramme de Gantt ci-dessous en donne la vue d'ensemble : les phases
d'analyse et de conception y précèdent le développement, mais elles ne s'y
substituent pas --- conformément à la démarche agile, une part de conception a
été refaite à chaque sprint, au vu de ce que la revue précédente avait révélé.

\figProjet[0.95]{gantt.png}{Diagramme de Gantt du projet}{gantt}

\section*{Conclusion}
\addcontentsline{toc}{section}{Conclusion}

Ce chapitre a permis de cerner le contexte du projet : un institut de beauté
dont l'activité, entièrement pilotée à la main, souffrait d'une information
unique, localisée et périssable. La solution retenue --- une API portant les
règles métier, une application mobile pour les clientes et un back-office pour
l'institut --- a été justifiée, ses enjeux et ses contraintes énoncés, et la
méthode de conduite du projet exposée. Le chapitre suivant traduit ces besoins
en une analyse et une conception détaillées.
